\documentclass[11pt]{article}
\usepackage{amsmath,amssymb,amsthm,hyperref}
\newtheorem{theorem}{Theorem}
\newtheorem{lemma}{Lemma}
\newtheorem{corollary}{Corollary}
\newcommand{\E}{\mathbb E}
\newcommand{\Poly}{\mathcal P}
\title{Positive completion of rational step data for two exceptional variables}
\author{Research proof: independently reviewed}
\date{30 September 2026}
\begin{document}\maketitle
\noindent\textbf{Status.} The complete proof below has been independently
reviewed by the root, size-bounds, and rational-design research agents.
\medskip

\noindent\textbf{Scope.} This note constructs two finite equal-weight real
random variables which, together with $m$ independent continuous uniform
variables, have all orders equally likely. It does not yet discretize the
remaining $m$ variables. In particular, it does not claim two polynomial-size
exceptional dice in a completely finite family.

\begin{theorem}[Positive completion]\label{completion}
Let $m\ge1$, $0<b_1<\cdots<b_K<1$, and suppose the equal-weight rule at
the $b_q$ integrates every polynomial of degree at most $m+1$ on $[0,1]$.
Write $b_0=0,b_{K+1}=1$ and
$\Delta=\min_{0\le q\le K}(b_{q+1}-b_q)$.
Suppose $0=r_0<r_1<\cdots<r_K<r_{K+1}=1$ satisfy
\begin{equation}\label{step}
 \frac1K\sum_{q=1}^K r_q b_q^j=\frac1{j+2},\qquad 0\le j\le m.
\end{equation}
There is an absolute $c>0$ such that, if
\begin{equation}\label{close}
 \max_q|r_q-b_q|\le c\frac{\Delta^{5/2}}{(m+1)^3},
\end{equation}
then a density $f$ on $[0,1]$, polynomial of degree at most $m$ on each
interval $(b_q,b_{q+1})$, has the following properties:
\begin{enumerate}
\item $1/2\le f\le3/2$ almost everywhere;
\item $\int_0^{b_q}f=r_q$;
\item if $A$ has density $f$, $B$ is uniform on the $K$ nodes $b_q$, and
$U_1,\ldots,U_m$ are independent uniform variables on $[0,1]$, then all
$(m+2)!$ strict orders of these independent variables have probability
$1/(m+2)!$.
\end{enumerate}
If all $r_q$ are rational with a common denominator $L$, then, after
multiplying $L$ by an integer if necessary, $A$ can instead be chosen
uniform on $L$ points in $\bigcup_q(b_q,b_{q+1})$, with the same conclusion.
It is enough that every integer $L(r_{q+1}-r_q)$ be at least $C(m+1)^2$.
\end{theorem}

Define the tail step functions
\[
 S_j(a)=\frac1K\sum_{b_q>a}b_q^j,\qquad 0\le j\le m,
\]
and the finite-dimensional space
\[
 V=\operatorname{span}\{a^j:0\le j\le m\}
   +\operatorname{span}\{a^rS_s(a):r,s\ge0,\ r+s\le m\}.
\]
Put $S=\operatorname{span}\{1,S_0,\ldots,S_m\}\subseteq V$.
All norms and orthogonal projections below use Lebesgue measure on $[0,1]$.
Let $P_0$ be orthogonal projection onto the space of functions constant
on each of the $K+1$ intervals $(b_q,b_{q+1})$.

\begin{lemma}[A stable quotient by the step functions]\label{quotient}
For every $F\in V$,
\[
 \operatorname{dist}_{L^2}(F,S)
 \le C(m+1)^2\Delta^{-1}\|(I-P_0)F\|_2.
\]
Moreover $V\cap\operatorname{ran}P_0=S$.
\end{lemma}
\begin{proof}
The jump of $F$ at $b_q$ is $-D(b_q)/K$ for a polynomial $D$ of degree
at most $m$. If $D(a)=\sum_jd_ja^j$, the function
$G=\sum_jd_jS_j\in S$ has the same jumps. Consequently $F-G$ is a
continuous piecewise polynomial on $[0,1]$. Adding a constant to $G$,
we may arrange $\int(F-G)=0$. The ordinary Poincare inequality gives
\[
 \|F-G\|_2\le C\|F'\|_2.
\]
On an interval $I$ of length at least $\Delta$, the polynomial inverse
Markov inequality gives
\[
 \|F'\|_{L^2(I)}
 \le C(m+1)^2|I|^{-1}
       \|F-\operatorname{mean}_I F\|_{L^2(I)}.
\]
Summing its square over the cells proves the claimed bound. If
$(I-P_0)F=0$, the bound gives $F\in S$, proving the intersection claim.
\end{proof}

\begin{proof}[Proof of Theorem~\ref{completion}: positive density]
On $(b_q,b_{q+1})$ set
\[
 g=\frac{r_{q+1}-r_q}{b_{q+1}-b_q}.
\]
Then $\int g=1$, $\int_0^{b_q}g=r_q$, and
\begin{equation}\label{gclose}
 \|g-1\|_\infty\le\frac{2\max_q|r_q-b_q|}{\Delta}.
\end{equation}
The step equations imply
\[
 \int gS_j=\frac1K\sum_q b_q^j r_q=\frac1{j+2}
 =\frac1K\sum_q b_q^{j+1}=\int S_j,
\]
so $g-1$ is orthogonal to $S$. On $W=(I-P_0)V$ define
\[
 \Lambda((I-P_0)F)=\int(g-1)F.
\]
Lemma~\ref{quotient} makes this well defined and gives
\[
 \|\Lambda\|\le C(m+1)^2\Delta^{-1}\|g-1\|_2.
\]
The Riesz representative $h\in W$ therefore satisfies
\[
 \int hF=\int(g-1)F\quad(F\in V),\qquad
 \|h\|_2\le C(m+1)^2\Delta^{-1}\|g-1\|_2.
\]
Here $h$ has zero integral on each cell and is polynomial of degree at
most $m$ there. The polynomial $L^2$--$L^\infty$ estimate on each cell
now yields
\[
 \|h\|_\infty\le C(m+1)\Delta^{-1/2}\|h\|_2
 \le C(m+1)^3\Delta^{-3/2}\|g-1\|_2.
\]
Set $f=g-h$. This retains every prescribed cell mass and satisfies
$\int fF=\int F$ for all $F\in V$. Equations~\eqref{close} and
\eqref{gclose}, with sufficiently small absolute $c$, give
$1/2\le f\le3/2$.

For completeness, these tests imply the full order assertion. For
$r+s\le m$,
\[
 \E[A^rB^s\mathbf1_{A<B}]
 =\int f(a)a^rS_s(a)\,da
 =\frac1K\sum_q\frac{b_q^{r+s+1}}{r+1}
 =\frac1{(r+1)(r+s+2)}.
\]
These are exactly the corresponding moments of two independent
continuous uniforms. Ordinary moments of $A$ and $B$ through degree $m$
are also correct. Subtracting the above triangle moments from their
product moments gives the corresponding identities for $B<A$.
For any prescribed order of $A,B,U_1,\ldots,U_m$, conditioning on
$A=a<B=b$ gives a polynomial proportional to
$a^i(b-a)^j(1-b)^k$, where $i+j+k=m$; the opposite triangle is identical
with the labels exchanged. Thus every such probability equals its
all-continuous uniform value $1/(m+2)!$.
\end{proof}

\begin{lemma}[Interior equal-weight quadrature]\label{interior}
Let a probability density $w$ on $[0,1]$ satisfy fixed bounds
$0<c_0\le w\le C_0$. For every integer $N\ge C_{c_0,C_0}(m+1)^2$
there is a degree-$m$ equal-weight quadrature for $w$ whose nodes all
lie in $(0,1)$.
\end{lemma}
\begin{proof}
Theorem~1.1 of Gilboa--Peled gives this bound with nodes in the closed
interval for any density bounded above and below by fixed positive
constants. We explain how to move the support strictly inside without
changing the degree-$m$ moments or these density bounds.

Take $J=[\tau,1-\tau]$, where $\tau=c(m+1)^{-2}$ and $c>0$ is sufficiently
small in terms of $c_0,C_0$. Let $K_m^J(t,x)$ be the reproducing kernel
for degree-$m$ polynomials in $L^2(J)$. The usual Legendre kernel bound
gives $|K_m^J(t,x)|\le C(m+1)^2$ for $x,t\in J$. The same bound, with
a changed absolute constant, holds for $x\in J,t\in[0,1]$: the affine
coordinate of $t$ is at most $1+C\tau$ in absolute value, and
$|P_j(1+v)|\le\exp(Cj\sqrt v)$ for $0\le v\le C\tau$.
Define on $J$
\[
 k(x)=\int_{[0,1]\setminus J}w(t)K_m^J(t,x)\,dt.
\]
Then $\|k\|_\infty\le C C_0\tau(m+1)^2\le c_0/2$ and, for every
polynomial $P$ of degree at most $m$,
\[
 \int_J P(x)k(x)\,dx
 =\int_{[0,1]\setminus J}P(t)w(t)\,dt.
\]
The density $\widetilde w=(w+k)\mathbf1_J$ has total mass one, agrees
with $w$ on these moments, and after affine rescaling of $J$ has
uniform positive lower and upper density bounds. Apply the closed
interval theorem on $J$. Its nodes, including possible endpoints of
$J$, belong to $(0,1)$ as required.
\end{proof}

\begin{proof}[Proof of Theorem~\ref{completion}: finite $A$]
The mass of cell $q$ is $\gamma_q=r_{q+1}-r_q=N_q/L$.
Rescaling the conditional density $f/\gamma_q$ from that cell to $[0,1]$
gives density bounds $1/3$ and $3$, since $f$ and its cell average both
lie in $[1/2,3/2]$. If $N_q\ge C(m+1)^2$, Lemma~\ref{interior} gives
a degree-$m$ equal-weight quadrature with exactly $N_q$ nodes strictly
inside the cell. Give all $\sum_qN_q=L$ nodes mass $1/L$. This preserves
the integral of every function which is polynomial of degree at most
$m$ in each cell, hence of every $F\in V$. The order argument above
applies unchanged, and no $A$ node coincides with a $B$ node.
\end{proof}

\begin{corollary}[A polynomial-size conditional conclusion]\label{poly}
Suppose, as $m$ varies, the step data in Theorem~\ref{completion} can
be chosen with $K$, $\Delta^{-1}$, and a common denominator $L$ all
bounded by fixed powers of $m+1$, and with~\eqref{close}. Then two
finite equal-weight independent variables with polynomially many
atoms each can coexist with $m$ independent continuous uniforms in
a completely order-fair family.
\end{corollary}
\begin{proof}
Every cell mass is at least $\Delta/2$. Multiplying $L$ by an integer
at most $1+C(m+1)^2/(L\Delta)$ makes all cell counts large enough.
The resulting total count remains polynomial, and the theorem applies.
\end{proof}

\begin{lemma}[An alternative with prescribed truncated moments]\label{prefix}
In the hypotheses of Theorem~\ref{completion}, put $d_q=r_q-b_q$,
including $d_0=d_{K+1}=0$. For any integer $D\ge m$, if
\[
 \max_q|d_q|\le c\Delta/(D+1)^2,
\]
there is a density $f$ between $1/2$ and $3/2$, polynomial of degree
at most $D$ on each cell, such that
\[
 \int_0^{b_q} a^s f(a)\,da
 =\frac{b_q^{s+1}}{s+1}+d_q b_q^s,
 \qquad 0\le s\le D.
\]
It yields the same hybrid order-fair family as in the theorem and
may be discretized cellwise in the same way.
\end{lemma}
\begin{proof}
For $I_q=[b_q,b_{q+1}]$ let $K_D^{I_q}$ be its degree-$D$ Lebesgue
reproducing kernel. On that cell set
\[
 f(a)=1+d_{q+1}K_D^{I_q}(b_{q+1},a)
          -d_qK_D^{I_q}(b_q,a).
\]
The kernel is bounded in absolute value by $(D+1)^2/|I_q|$ on the
closed cell. This proves positivity and, after telescoping the cell
moment corrections, the asserted truncated moments. In particular
the cell masses are the prescribed rational differences of ranks.
The full moments of $A$ are uniform because $d_{K+1}=0$. For $r+s\le m$,
\[
 \E[A^rB^s\mathbf1_{A<B}]
 =\frac1K\sum_q\left(\frac{b_q^{r+s+1}}{r+1}
                          +d_qb_q^{r+s}\right)
 =\frac1{(r+1)(r+s+2)},
\]
since~\eqref{step} makes $K^{-1}\sum_qd_qb_q^j=0$ for $j\le m$.
The triangle argument already proved finishes the claim.
\end{proof}

\paragraph{Input and remaining gap.}
The companion construction by the rational-design research agent
produces the step data with polynomial separation and denominator,
using a small Legendre perturbation followed by a quantitatively
conditioned implicit-function argument. This note isolates the
positive completion so that its conditioning can be independently
checked. An additional argument is required to replace the $m$
continuous uniforms by finite dice while retaining the two small atom
counts. The one-exceptional construction cannot simply be invoked:
two distinguished cut points permit surviving two-bump interaction terms.
That additional argument is now supplied by the independently reviewed
\path{../size_bounds/two_polynomial_exceptional_dice.tex} and its
clustered lifting companion. It uses rational box smoothing and preserves
all mixed boundary terms explicitly; the scope of the present theorem
itself remains the stated hybrid model.
Lemma~\ref{prefix} is a later elementary strengthening, independently
checkable from the reproducing-kernel identity; its special truncated
moment formula is not needed for the quotient argument above.

\begin{thebibliography}{1}
\bibitem{GP} S. Gilboa and R. Peled,
\emph{Chebyshev-type quadratures for doubling weights},
Constructive Approximation 45 (2017), 193--216, Theorem~1.1.
\href{https://arxiv.org/abs/1507.01505}{arXiv:1507.01505}.
\end{thebibliography}
\end{document}
